\section{Lock Freedom}
\label{sec:lock-freedom}

% PAR: section overview, the shaded parts of the earlier figures (conv:lf), no new rules, thm:lf mirrors thm:df | example: - | figure/rule: conv:lf, thm:lf | map: ty-ann

\subsection{Three Ways to Wait Forever}
\label{sec:lf-why}

% PAR: lock freedom as defined in def:lf, every pending thread eventually takes part | example: - | figure/rule: def:lf | map: sem-props
% PAR: order alone gives only deadlock freedom, a thread waits forever in three ways, each excluded by one rule on counters (fig:lf-failures) | example: rx-postpone, rx-delegate, rx-relink | figure/rule: fig:lf-failures | map: sem-props, t-lf-counters
% PAR: endless postponement, a thread holds an end across a loop that never acts on it, every action lowers the counter of the held end, a cost lowers its credit | example: rx-postpone | figure/rule: ex:postpone, fig:lf-failures, app:t-agecap | map: t-agecap

\begin{example}[Postponement \rocq{rx_postpone_rejected}]
\label{ex:postpone}
% Generated from map example rx-postpone (Postponement, a holding thread that never serves). Do not edit.
\leavevmode\par\nobreak\noindent
\begin{minipage}{\linewidth}
\begin{maplisting}
chat = !(f = λa. \lstkw{let} a = \lstkw{send} ⟨∗, a⟩ \lstkw{in} (\lstkw{force} f) a)
loop = !(f = λ⟨a, b⟩. \lstkw{let} ⟨u, a⟩ = \lstkw{recv} a \lstkw{in} \lstkw{let} ∗ = u \lstkw{in} (\lstkw{force} f) ⟨a, b⟩)
main = \lstkw{let} ⟨a, a′⟩ = \lstkw{new} \lstkw{in} \lstkw{let} ⟨b, b′⟩ = \lstkw{new} \lstkw{in}
       \lstkw{let} ∗ = \lstkw{spawn} (λ_. (\lstkw{force} chat) a′) \lstkw{in}
       \lstkw{let} ∗ = \lstkw{spawn} (λ_. \lstkw{close} b′) \lstkw{in}
       (\lstkw{force} loop) ⟨a, b⟩            \lstcmt{-- holds b forever, close b′ waits forever}
\end{maplisting}
\end{minipage}

% PAR: deadlock free, not lock free, the credit of b drops by the cost of each round | example: rx-postpone | figure/rule: app:t-agecap | map: t-agecap
\end{example}

% PAR: endless passing, an end sent back and forth, a send lowers the counter of each sent end by one | example: rx-delegate | figure/rule: ex:delegate, E-Send | map: r-act
\begin{example}[Endless delegation \rocq{rx_delegate_rejected}]
\label{ex:delegate}
% Generated from map example rx-delegate (Endless delegation, an end that travels forever). Do not edit.
\leavevmode\par\nobreak\noindent
\begin{minipage}{\linewidth}
\begin{maplisting}
pass = !(f = λ⟨e, x, y⟩. \lstkw{let} x = \lstkw{send} ⟨e, x⟩ \lstkw{in} \lstkw{let} ⟨e, y⟩ = \lstkw{recv} y \lstkw{in}
               (\lstkw{force} f) ⟨e, x, y⟩)
back = !(f = λ⟨x, y⟩. \lstkw{let} ⟨e, x⟩ = \lstkw{recv} x \lstkw{in} \lstkw{let} y = \lstkw{send} ⟨e, y⟩ \lstkw{in}
               (\lstkw{force} f) ⟨x, y⟩)
main = \lstkw{let} ⟨e, e′⟩ = \lstkw{new} \lstkw{in} \lstkw{let} ⟨x, x′⟩ = \lstkw{new} \lstkw{in} \lstkw{let} ⟨y, y′⟩ = \lstkw{new} \lstkw{in}
       \lstkw{let} ∗ = \lstkw{spawn} (λ_. \lstkw{close} e′) \lstkw{in}       \lstcmt{-- waits for e forever}
       \lstkw{let} ∗ = \lstkw{spawn} (λ_. (\lstkw{force} back) ⟨x′, y′⟩) \lstkw{in}
       (\lstkw{force} pass) ⟨e, x, y⟩                \lstcmt{-- e goes back and forth forever}
\end{maplisting}
\end{minipage}

% PAR: payload premise of E-Send rejects | example: rx-delegate | figure/rule: E-Send | map: r-act
\end{example}

% Listing in the supplement (rx-relink).
% PAR: endless links, each link hands the waiting end a new dual end, the new dual end has a lower counter, so a wait meets at most N links | example: rx-relink | figure/rule: fig:lf-failures, E-Link | map: r-act
% Three ways a thread of a deadlock-free program waits forever, and the rule
% on counters that excludes each. From map examples rx-postpone
% (rx_postpone_rejected), rx-delegate (rx_delegate_rejected) and rx-relink
% (rx_relink_rejected, listing in the supplement), blocks t-lf-counters,
% t-agecap, r-act. The run-time column is def:counters, the typing column
% the premises of fig:typing-core and app:t-agecap.
\begin{figure}[tb]
  \small
  \centering
  \begin{tabular}{@{}p{0.25\linewidth}@{\quad}p{0.32\linewidth}@{\quad}p{0.33\linewidth}@{}}
    \toprule
    the thread holding $\dual{o}$ & at run time & in typing \\
    \midrule
    acts on other ends forever (\cref{ex:postpone}) &
    each action elsewhere lowers $f(\dual{o})$ by one &
    a term of cost $\varepsilon$ asks $\varepsilon \leq a$ of the credit $a$ of $\dual{o}$ and lowers it to $a - \varepsilon$, and a force costs $1$ \\[2pt]
    sends $\dual{o}$ on forever (\cref{ex:delegate}) &
    a send lowers the counter of each sent end by one &
    \ruleref{E-Send} asks $\ty' - 1 \leq \payload{\ty}{x}$ of the payload \\[2pt]
    links $\dual{o}$ forever, so $o$ meets a new dual end each time &
    the new dual end has a lower counter than the old one, so $o$ meets at most $N$ links &
    \ruleref{E-Link} asks $a_1 + m_2 > N$ and $a_2 + m_1 > N$ \\
    \bottomrule
  \end{tabular}
  \Description{A table with three rows, one per way in which a thread waits forever on an end o, with the run-time counter that excludes it and the typing premise.}
  \caption{Three ways a thread waiting on $o$ waits forever in a deadlock-free program, and the counters that exclude each. The third row is the program of rejected relinking in the supplement.}
  \label{fig:lf-failures}
\end{figure}


% PAR: counters alone do not exclude cyclic waits, orders and counters in one judgment | example: ex-milner | figure/rule: def:lf-ann | map: t-lf-counters, fu-vs
% GAP: no map example for "counters alone admit a cyclic wait"

\subsection{Counters and Costs}
\label{sec:lf-annotations}

% PAR: LF typing = DF typing + annotations, the shaded parts of every figure, two kinds of counting | example: - | figure/rule: conv:lf, def:lf-ann | map: ty-grammar, t-judgment, ty-ann
\begin{definition}[Counters and costs \rocq{tsum_ok} \rocq{tage}]
\label{def:lf-ann}
Lock-free typing counts in two ways.
\begin{itemize}
  \item Counters. At run time each held end $o$ has a counter $f(o)$, the number of actions its thread may still take on other ends before it acts on $o$. An action on $o$ resets $f(o)$ to the reset value of $o$. Every action of the thread lowers the counters of its other held ends by one, and a send lowers the counter of each end it sends by one.
  \item Cost. The cost of a term bounds the number of actions it takes before it returns.
\end{itemize}
The types record both.
\begin{center}
\small
\begin{tabular}{@{}l@{\quad}p{0.68\linewidth}@{}}
  $o : \sty_x^{a,m}$ & credit $a$, a static lower bound $a \leq f(o)$, and tickets $\tk(o) = m$, the reset value, so $f(o) \leq m$ \\
  $N$ & bound on the tickets of the two ends of a channel, $\tk(o) + \tk(\dual{o}) \leq N$ \\
  $(\ty \arr{u}{l}[c]\, \tyB)^{a}$ & credit $a \leq f(o)$ for each captured end $o$ that is not idle, latent cost $c$ of the body \\
  $\tm : \ty\ !\ \varepsilon$ & cost $\varepsilon$ of $\tm$ \\
  $\lenv - \varepsilon$ & $\lenv$ with each credit lowered by $\varepsilon$, defined when $\varepsilon$ is at most each credit of a closure and of an end that is not idle (\cref{app:t-agecap})
\end{tabular}
\end{center}
An end is idle when its session type unfolds to $\Skip$. Every annotation is in $\mathbb{N} \cup \{\infty\}$, and $\infty - c = \infty$.
\end{definition}

% PAR: worked example, the cycle at N = 4, two tickets per end, closing c costs one, the credit of d falls from 2 to 1 and so does its counter | example: ex-cycle | figure/rule: fig:lf-example | map: t-lf-counters, t-agecap
\begin{example}[Cycle at $N = 4$ \rocq{cy_lock_free}]
\label{ex:cycle}
% Generated from map example ex-cycle (Cycle). Do not edit.
\leavevmode\par\nobreak
\begin{maplisting}
cycle = let ⟨c, c′⟩ = new in          -- c at ρ
        let ⟨d, d′⟩ = new in          -- d at ρ·0, C = {ρ ≺ (ρ·0)‾, ρ‾ ≺ ρ·0}
        let ∗ = spawn (λu. let ∗ = u in let ∗ = close c′ in close d′) in
        let ∗ = close c in close d
\end{maplisting}

\end{example}
% Worked example of the lock-free annotations: the main thread of map
% example ex-cycle at N = 4 (theories/examples/LFCycle.v, cy_typedN with
% e_new' and e_new_at' of examples/LFTools.v: tickets N/2 = 2 per end).
% Types from E-New, E-LetU, E-Close and Age-Credit; counters from the
% fuel obligations of an action (def:fuel).
\begin{figure}[tb]
  \small
  \centering
  \begin{tabular}{@{}l@{\qquad}l@{\qquad}l@{\qquad}c@{\qquad}c@{}}
    \toprule
    after the line & $c$ & $d$ & cost & $f(d)$ \\
    \midrule
    $\tletp{c}{c'}{\tnew}{}$ & $\Close_{\rho}^{2,2}$ & & $0$ & \\
    $\tletp{d}{d'}{\tnew}{}$ & $\Close_{\rho}^{2,2}$ & $\Close_{\fld{\rho}{0}}^{2,2}$ & $0$ & $2$ \\
    $\tletu{\tspawn\,(\ldots)}{}$ & $\Close_{\rho}^{2,2}$ & $\Close_{\fld{\rho}{0}}^{2,2}$ & $0$ & $2$ \\
    $\tletu{\tclose\,c}{}$ & & $\Close_{\fld{\rho}{0}}^{1,2}$ & $1$ & $1$ \\
    $\tclose\,d$ & & & $1$ & \\
    \bottomrule
  \end{tabular}
  \Description{A table with one row per line of the main thread of the cycle, giving the annotated types of the ends c and d, the cost of the line and the counter of d.}
  \caption{The main thread of \cref{ex:cycle} at $N = 4$. Each row gives the types of $c$ and $d$ after the line, the cost of the line and the counter $f(d)$ at run time, so closing $c$ ages the credit of $d$ from $2$ to $1$ and lowers $f(d)$ from $2$ to $1$.}
  \label{fig:lf-example}
\end{figure}


% PAR: shaded parts of E-App, E-Pair, E-Case, sequencing of costs, the held closure asks the cost of the argument below its credit | example: - | figure/rule: fig:typing-core | map: r-fun, r-data
% PAR: actions reset to m, cost one | example: ex-cycle | figure/rule: E-Send, E-Close | map: r-act
% PAR: payload credit, a sent end loses one | example: rx-delegate | figure/rule: E-Send | map: r-act
% PAR: link tickets, a1 + m2 > N | example: rx-relink | figure/rule: E-Link | map: r-act
% PAR: creation splits the tickets, m1 + m2 <= N | example: ex-cycle | figure/rule: E-New | map: r-new
% PAR: loops charged one force, loops that hold an end | example: ex-guard-neg, ex-guard-pos | figure/rule: E-Force | map: r-boxes, t-agecap
% PAR: process typing at N, credits and tickets of interface ends in Sigma | example: - | figure/rule: fig:process-typing | map: p-typing
% PAR: erasure gives DF typing, conservativity | example: - | figure/rule: thm:conserv | map: th-conserv

\begin{theorem}[Conservativity \rocq{conservative}]
\label{thm:conserv}
Let every credit, cost and ticket in $\lenv$ and $\ty$ be $\infty$. Then $\cGamma ; \lenv \vdash^{F}_{W,\infty} \tm : \ty\ !\ \varepsilon$ for some $\varepsilon$ if and only if $\cGamma ; \lenv \vdash^{F}_{W} \tm : \ty$.
\end{theorem}

% Listing in the supplement (ex-controller).
% PAR: idle end held across an unbounded loop, its credit truncated at 0 by the cost | example: ex-controller | figure/rule: app:t-agecap | map: t-agecap

\subsection{Ranked Worlds}
\label{sec:lf-ranked}

% PAR: consistent worlds become ranked worlds, ordinal ranks, well-founded | example: ex-milner | figure/rule: def:ranked, fig:world-example | map: w-world, w-wf, t-lf-wf
\begin{definition}[Ranked worlds \rocq{wranked}]
\label{def:ranked}
\begin{gather*}
  \ranked\,W \iff \exists \varrho : \Point \to \mathsf{Ord}.\ \forall x, y.\ x \ordW y \Rightarrow \varrho(x) < \varrho(y)
\end{gather*}
$\ranked\,W$ implies $\cons\,W$ (\cref{app:w-wf}). At a finite $N$, \ruleref{P-Term}, \ruleref{P-Connect} and the lock-free judgment read $\ranked$ where \cref{sec:typing,sec:relation} read $\cons$ (\cref{conv:lf}), and \ruleref{E-New} asks $\rankable$ (\cref{def:creation}).
\end{definition}
% PAR: the scheduler's world is ranked, ranks read the field label first | example: ex-milner | figure/rule: fig:world-example | map: t-lf-wf
% PAR: they differ only on infinite worlds | example: ex-forkjoin | figure/rule: def:ranked | map: w-wf
% PAR: new_ok_N, the new tree rankable | example: ex-gennew | figure/rule: E-New | map: r-new
% PAR: recursive fork-join, infinite descending chain, typable in DF, no LF typing | example: ex-forkjoin | figure/rule: ex:forkjoin | map: r-new, t-lemmas
\begin{example}[Recursive fork-join with replies \rocq{fjr_deadlock_free}]
\label{ex:forkjoin}
% Generated from map example ex-forkjoin (Recursive fork-join with replies). Do not edit.
\leavevmode\par\nobreak\noindent
\begin{minipage}{\linewidth}
\begin{maplisting}
main = \lstkw{let} ⟨w, w′⟩ = \lstkw{new} \lstkw{in}                      \lstcmt{-- w at ρ, facts H at \ensuremath{\bar{ρ}}}
       \lstkw{let} ∗ = \lstkw{spawn} (λu. (\lstkw{force} wk) w′) \lstkw{in} \lstkw{close} w
wk   = !(f = λw. \lstkw{let} ⟨c₁, c₁′⟩ = \lstkw{new} \lstkw{in} \lstkw{let} ⟨c₂, c₂′⟩ = \lstkw{new} \lstkw{in}   \lstcmt{-- at ι·0 and ι·1}
             \lstkw{let} ∗ = \lstkw{spawn} (λu. (\lstkw{force} f) c₁′) \lstkw{in}
             \lstkw{let} ∗ = \lstkw{spawn} (λu. (\lstkw{force} f) c₂′) \lstkw{in}
             \lstkw{let} ∗ = \lstkw{close} c₁ \lstkw{in} \lstkw{let} ∗ = \lstkw{close} c₂ \lstkw{in} \lstkw{close} w)    \lstcmt{-- reply after both}
\end{maplisting}
\end{minipage}

% PAR: infinite descending chain, consistent, deadlock free, typable in the deadlock-free instance | example: ex-forkjoin | figure/rule: E-New | map: r-new, w-wf
\end{example}
% PAR: a waiting chain that grows forever | example: rx-chain | figure/rule: E-New | map: r-new

% Listing in the supplement (rx-chain).
% PAR: ranks admit no infinite descent, rejected at a finite bound | example: rx-chain | figure/rule: def:ranked | map: t-lf-wf

\subsection{The Lock-Free Relation}
\label{sec:lf-relation}

% PAR: cost table kappa_L, the same E | example: - | figure/rule: fig:term-rel | map: ct-costs, e-rel
% PAR: the relation reads the counters of def:lf-ann at run time, E extended with counters, the cost bound is the cost, no world, no readiness | example: ex-cycle | figure/rule: def:counters | map: fu, fu-vs
% PAR: mechanization insight, E and the counters kept as two fixed points, not merged | example: - | figure/rule: def:counters | map: fu-vs
\begin{definition}[Counters of a term \rocq{FuelE}]
\label{def:counters}
A term keeps its counters, $\tm \in \Fuel{N}{\ix}{\cPhi}{\sigma\ !\ d}$, with the components
\begin{center}
\small
\begin{tabular}{@{}l@{\quad}l@{}}
  $\cPhi$ & held ends, each with its protocol $\sty$, counter $f$ and tickets $m$, written $\sty^{f,m}$ \\
  $\sigma$ & type of the result, or $\mathsf{thread}$ for a thread, which returns holding no end \\
  $d$ & cost bound, the number of actions $\tm$ may take before it returns
\end{tabular}
\end{center}
It reads no point and no world and has no readiness clause. It holds at $\ix$ when the following hold (\cref{app:fu,app:fu-obls,app:fv}).
\begin{itemize}
  \item (stop) A value has its counters at least its credits. A closure has the counter of each captured end that is not idle at least its credit, and a box holds no end.
  \item (step) Each step $\tm \xrightarrow{\ell} \tm'$ changes the held ends as its obligation in \cref{fig:term-rel} does, and $\tm'$ keeps its counters at each $\ixs < \ix$ as follows.
  \item An action on $o$ asks $d \geq 1$ and a counter of at least $1$ on each other held end that is not idle. It lowers $d$ and those counters by one and resets the counter of $o$ to its tickets, or drops $o$ at a close.
  \item A force asks $d \geq 1$ and a counter of at least $1$ on each held end that is not idle, and lowers $d$ and every counter by one.
  \item A send passes on the ends of its payload, each with a counter of at least $1$, and lowers their counters by one. A creation starts its two ends at tickets $m_1$ and $m_2$ with $m_1 + m_2 \leq N$, each counter at its tickets. A link of two ends with counters $f_1, f_2$ and tickets $m_1, m_2$ asks $f_1 + m_2 > N$ and $f_2 + m_1 > N$. A $\beta$ step, a spawn and a link lower no counter.
\end{itemize}
\end{definition}
% PAR: counter obligations, force lowers every counter, an action needs counters of at least one | example: ex-guard-neg | figure/rule: def:counters | map: fu-obls
% PAR: link, the counters of the new dual ends fall | example: ex-link | figure/rule: app:fu-obls | map: fu-obls
% PAR: received closures own their ends, exactness | example: ex-forwarder | figure/rule: app:fv | map: fv

% PAR: LF process relation = E at kappa_L plus the counters under one ownership map, the term relation extended with counters at cost bound infinity | example: ex-cycle | figure/rule: def:lf-rel | map: soup
\begin{definition}[Lock-free process relation \rocq{SoupK}]
\label{def:lf-rel}
An annotated process is $a = (\proc_a, K, W_a)$ with
\begin{center}
\small
\begin{tabular}{@{}l@{\quad}l@{}}
  $\proc_a$ & process \\
  $K$ & ownership map, $K(o) = (t, x, f, m)$ the owner, point, counter and tickets of each held end $o$ \\
  $W_a$ & ranked world of the run so far
\end{tabular}
\end{center}
$a \in \LFrel{N,W}{\ix}{\senv}{\ty}$ holds when $W \sqsubseteq W_a$, $\tfree\,\ty$, $K$ fits $\proc_a$ and $\senv$ (\cref{app:soup}), and for some session type and protocol per held end each thread $\tm = \proc_a(t)$ has
\[
  \tm \in E[\kappa_L]^{W_a}_{\ix}\llbracket \cTheta_t \vdash \tyB_t \rrbracket^{\varnothing}
  \qquad\text{and}\qquad
  \tm \in \Fuel{N}{\ix}{\cPhi_t}{\mathsf{thread}\ !\ \infty}
\]
\begin{itemize}
  \item $\cTheta_t$, $\cPhi_t$: the ends $K$ gives $t$, at their points and at their counters and tickets.
  \item $\tyB_t$: $\ty$ for the main thread and $\ast$ otherwise.
  \item $E[\kappa_L]$: the term relation of \cref{fig:term-rel} under the cost table $\kappa_L$, which charges $1$ for $\tforce$ and each session action and $0$ for $\beta$, $\tnew$, $\tspawn$ and $\tlink$ (\cref{app:ct-costs}).
  \item $K$ fits $\proc_a$ and $\senv$: each held end belongs to a thread of $\proc_a$ and is restricted in $\proc_a$ or listed in $\senv$, both ends of a restricted channel are held at dual points and dual protocols with tickets of sum at most $N$, each counter is at most its tickets and each ticket count at most $N$, and each end of $\senv$ carries the point and protocol of its entry of $\senv$ and at most the tickets of that entry.
\end{itemize}
\end{definition}

% PAR: lock-free judgment over ranked extensions | example: - | figure/rule: def:lf-judgment | map: sj-all
\begin{definition}[Lock-free judgment \rocq{lf_sems}]
\label{def:lf-judgment}
\[
  \senv \semlf^{N}_{W} \proc : \ty \iff \forall W' \sqsupseteq W.\ \ranked\,W' \Rightarrow \forall \ix.\ \exists a = (\proc_a, K, W_a).\ \proc_a \sim \proc \land a \in \LFrel{N,W'}{\ix}{\senv}{\ty}
\]
Here $\proc_a \sim \proc$ holds when $\proc_a$ is $\proc$ with its threads renamed (\cref{app:syn-procs}).
\end{definition}

% PAR: lock-free composition of two components at own worlds | example: ex-cs-split | figure/rule: thm:lf-compose | map: th-compose, w-ext
\begin{theorem}[Lock-free composition \rocq{lf_p_term} \rocq{lf_p_connect}]
\label{thm:lf-compose}
Let $N \in \mathbb{N}$.
\begin{enumerate}
  \item If the premises of \ruleref{P-Term} at $N$ hold apart from $\ranked\,W$, then $\senv \semlf^{N}_{W} (C_0, \{\tm\}) : \ty$.
  \item If $\pfree\,\senv_1$, $\pfree\,\senv_2$, $\pfree\,\senv_C$, $\senv_1 \cup \senv_C \semlf^{N}_{W_1} \proc_1 : \ty$ and $\senv_2 \cup \dual{\senv_C} \semlf^{N}_{W_2} \proc_2 : \tyunit$ and $\tk(o) + \tk(\dual{o}) \leq N$ for each $o \in \mathrm{dom}\,\senv_C$, then $\senv_1 \cup \senv_2 \semlf^{N}_{W_1 \wjoin W_2} \proc_1 \cpar{\senv_C} \proc_2 : \ty$.
\end{enumerate}
\end{theorem}

% PAR: shaded lines of thm:fundamental-terms, thm:fundamental-proc, thm:adequacy | example: ex-cs | figure/rule: sec:fundamental | map: th-terms, th-fund-proc, th-adeq
\begin{theorem}[Lock-free fundamental theorems and adequacy \rocq{live_fundamental_ranked} \rocq{lf_fundamental} \rocq{lf_sems_adequacy}]
\label{thm:lf-fundamental}
Let $N \in \mathbb{N}$.
\begin{enumerate}
  \item If $\varnothing ; \lenv \vdash^{F}_{W,N} \tm : \ty\ !\ \varepsilon$, then $|\lenv| \semt[\kappa_L]^{F}_{W'} \tm : |\ty|$ for each consistent $W' \sqsupseteq W$, where $|{\cdot}|$ erases the annotations and $\semt[\kappa_L]$ is the term judgment of \cref{def:sem-judgments} over $E[\kappa_L]$.
  \item If $\senv \vdash^{N}_{W} \proc : \ty$, then $\senv \semlf^{N}_{W} \proc : \ty$.
  \item If $\ranked\,W$ and $\varnothing \semlf^{N}_{W} \proc : \ty$, then $\proc$ is lock free.
\end{enumerate}
\end{theorem}
% PAR: adequacy measure, rank of the awaited point plus counter | example: - | figure/rule: def:ranked, def:lf-rel | map: t-lf-wf, th-adeq

\subsection{Lock Freedom of Well-Typed Programs}
\label{sec:lf-theorem}

% PAR: the main theorem, mirror of thm:df, corollary of thm:lf-fundamental (2) and (3) | example: ex-cycle | figure/rule: thm:lf, thm:df | map: th-sound
\begin{theorem}[Lock freedom \rocq{lf_soundness}]
\label{thm:lf}
Every closed process that is well typed at a finite bound is lock free. Assume excluded middle and functional choice. If $N \in \mathbb{N}$ and $\varnothing \vdash^{N}_{W} \proc : \ty$, then $\proc$ is lock free (\cref{def:lf}).
\end{theorem}
% PAR: footnote slot, excluded middle and functional choice enter only here and in thm:lf-fundamental, deadlock freedom needs neither | example: - | figure/rule: thm:lf, thm:df | map: sr-ord, th-sound
% PAR: the cycle is lock free at N = 4 | example: ex-cycle | figure/rule: ex:cycle, thm:lf | map: th-sound

% Listing in the appendix (app:ex-cs-split), cut from the body for space.
% PAR: worlds of their own, join ranked, lock-free composition | example: ex-cs-split | figure/rule: app:ex-cs-split, thm:lf-compose | map: th-compose

\FloatBarrier
